%\documentclass[twocolumn]{article}\
\documentclass{article}
\usepackage{amsmath, amssymb, cancel, mathtools, bm}
\usepackage[left=.5in, right=.5in, top=1in, bottom=1in]{geometry}
\usepackage[most]{tcolorbox}
\usepackage[skip=1em,indent=0pt]{parskip}
\setlength{\parindent}{0pt}
\newcommand{\statvec}[1]{\underset{\sim}{\bm{#1}}} % Vector symbol (tilde under vector)
\DeclareMathOperator{\EX}{\mathbb{E}} % Expected Value symbol
\DeclareMathOperator{\ext}{\EX_{\theta}} % Expected Value symbol
\DeclareMathOperator{\vart}{Var_{\theta}} % Expected Value symbol
\newcommand{\indep}{\perp\!\!\!\!\perp} % Independence symbol
\newcommand{\real}{\mathbb{R}} % Simplified 'Reals' indicator
\newcommand{\pto}{\overset{P}{\to}}
\newcommand{\asto}{\overset{a.s.}{\to}}
\newcommand{\rthto}{\overset{L^r}{\to}}
\newcommand{\dto}{\overset{D}{\to}}
\newcommand{\xtb}{x^{\top}\statvec{\beta}}
\newcommand{\xitb}{x_i^{\top}\statvec{\beta}}
% Force all aggregate symbols to always put values above/below
\let\Oldint=\int
\let\Oldsum=\sum
\let\Oldprod=\prod
\let\Oldbigcup=\bigcup
\let\Oldbigcap=\bigcap
\let\Oldlim=\lim
\renewcommand{\int}{\Oldint\limits} 
\renewcommand{\sum}{\Oldsum\limits} 
\renewcommand{\prod}{\Oldprod\limits} 
\renewcommand{\bigcup}{\Oldbigcup\limits} 
\renewcommand{\bigcap}{\Oldbigcap\limits} 
\renewcommand{\lim}{\Oldlim\limits} 
\newcommand{\infint}{\int_{-\infty}^{\infty}}
\newcommand{\iiddist}{\overset{\mathrm{iid}}{\sim}}
\newcommand{\norm}[1]{\left\lVert#1\right\rVert}
\newcommand{\boldred}[1]{\textbf{\textcolor{red}{#1}}}
\usepackage[linktoc=all]{hyperref}

\author{RJ Cass}
\title{STAT 637 - HW 2.2}

\begin{document}

\maketitle

%%%%%%%%%%%%%%%%%%%%%%%%%%%%%%%%%%%%%%%%%%%%%%%%%%%%%%%%%%%%%%%%%%%%%%%%%%%

\section{}

Show that the following distributions belong to the exponential family. For each distribution, identify: a natural parameter, a sufficient statistic, the mean, and the variance using the equations shown in class

Exponential Family: $f(y; \theta) = a(\theta)b(y)exp\{\sum_{j=1}^k w_j(\theta) t_j(y) \}$

Where $c(\theta) = ln(a(\theta))$:

Mean: $\EX(t(y)) = -\frac{c'(\theta)}{w'(\theta)}$

Var: $Var(t(y)) = \frac{w''(\theta)c'(\theta) - w'(\theta)c''(\theta)}{(w'(\theta))^3}$

(a) Gamma with unknown rate $\beta$, fixed $\alpha$: $f(y; \beta) = \frac{\beta^{\alpha}}{\Gamma(\alpha)}y^{\alpha-1}e^{-y\beta}$

Exponential Family: 
\begin{itemize}
    \item $a(\beta) = \frac{\beta^{\alpha}}{\Gamma(\alpha)}$
    \item $b(y) = y^{\alpha-1}$
    \item $w_1(\beta) = \beta$
    \item $t_1(y) = -y$
\end{itemize}

Natural Parameter: $\eta = \beta$

Sufficient statistic: $T(y) = \sum y$

Mean: 

\begin{quote}
$c'(\beta) = \frac{d}{d \beta} ln(\frac{\beta^{\alpha}}{\Gamma(\alpha)}) = \frac{d}{d \beta} \alpha ln(\beta) - ln(\Gamma(\alpha)) = \frac{\alpha}{\beta}$

$w'(\beta) = \frac{d}{d \beta} \beta = 1$

$\EX[y] = -(-1)\frac{\frac{\alpha}{\beta}}{1} = \frac{\alpha}{\beta}$
\end{quote}

Variance: 

\begin{quote}
$c''(\beta) = \frac{d}{d \beta} \frac{\alpha}{\beta} = -\frac{\alpha}{\beta^2}$

$w''(\beta) = \frac{d}{d \beta} 1 = 0$

$Var(y) = \frac{0 - 1(-\frac{\alpha}{\beta^2})}{1^3} = \frac{\alpha}{\beta^2}$
\end{quote}

(b) Poisson: $f(y; \lambda) = \frac{\lambda^y}{y!}e^{-\lambda}$

$\lambda^y = e^{y ln(\lambda)} \implies f(y; \lambda) = \frac{1}{y!}e^{-\lambda + yln(\lambda)}$

Exponential Family: 
\begin{itemize}
    \item $a(\lambda) = 1$
    \item $b(y) = \frac{1}{y!}$
    \item $w_1(\lambda) = \lambda, w_2(\lambda) = ln(\lambda)$
    \item $t_1(y) = -1, t_2(y) = y$
\end{itemize}

Natural Parameter: $\eta = (\lambda, ln(\lambda))$

Sufficient statistic: $T(y) = \sum y$

Mean: 

\begin{quote}
$c'(\lambda) = \frac{d}{d \lambda} ln(1) = 1$

$w'(\lambda) = \frac{d}{d \lambda} ln(\lambda) = \frac{1}{\lambda}$

$\EX[y] = \frac{1}{\frac{1}{\lambda}} = \lambda$
\end{quote}

Variance: 

\begin{quote}
$c''(\lambda) = \frac{d}{d \lambda} 1 = 0$

$w''(\lambda) = \frac{d}{d \lambda} \frac{1}{\lambda} = -\frac{1}{\lambda^2}$

$Var(y) = \frac{-\frac{1}{\lambda^2} - 0}{(-\frac{1}{\lambda})^3} = \frac{\frac{1}{\lambda^2}}{\frac{1}{\lambda^3}} = \lambda$
\end{quote}

(c) Negative Binomial with fixed $r$: $f(y; \pi) = \binom{y+r-1}{r-1} \pi^r (1-\pi)^y$

$(1-\pi)^y = e^{yln(1-\pi)} \implies f(y; \pi) = \binom{y+r-1}{r-1} \pi^r e^{yln(1-\pi)}$

Exponential Family: 
\begin{itemize}
    \item $a(\pi) = \pi^r$
    \item $b(y) = \binom{y+r-1}{r-1}$
    \item $w_1(\pi) = ln(1-\pi)$
    \item $t_1(y) = y$
\end{itemize}

Natural Parameter: $\eta = ln(1-\pi)$

Sufficient statistic: $T(y) = \sum y$

Mean: 

\begin{quote}
$c'(\pi) = \frac{d}{d \pi} ln(\pi^r) = \frac{d}{d \pi} r ln(\pi) = \frac{r}{\pi}$

$w'(\pi) = \frac{d}{d \pi} ln(1-\pi) = -\frac{1}{1-\pi}$

$\EX[y] = -\frac{\frac{r}{\pi}}{-\frac{1}{1-\pi}} = \frac{r (1-\pi)}{\pi}$
\end{quote}

Variance: 

\begin{quote}
$c''(\pi) = \frac{d}{d \pi} \frac{r}{\pi} = -\frac{r}{\pi^2}$

$w''(\pi) = \frac{d}{d \pi} -\frac{1}{1-\pi} = \frac{1}{(1-\pi)^2}$

$Var(y) = \frac{\frac{1}{(1-\pi)^2}\frac{r}{\pi} - \frac{r}{\pi^2}\frac{1}{1-\pi}}{(\frac{1}{1-\pi})^3} = \frac{r(1-\pi)^3}{\pi^2 (1-\pi)^2} = \frac{r(1-\pi)}{\pi^2}$
\end{quote}

%%%%%%%%%%%%%%%%%%%%%%%%%%%%%%%%%%%%%%%%%%%%%%%%%%%%%%%%%%%%%%%%%%%%%%%%%%%

\section{}

Poisson: $Var(Y) = \EX[Y]$

Binomial: $Var(Y) = \EX[Y](1-p)$

Negative Binomial: $Var(Y) = \frac{\EX[Y]}{p}$

(a) These models differ in a ratio of the 'success' parameter. In Poisson, the mean equals the variance, but in Binomial the variance is always smaller than the mean, and Negative Binomial the variance is larger. 

(b) Poisson may be inappropriate in samples where $\EX[Y]$ is very large, because this would also result in a very large variance. 

(c) If the variability is larger than the mean, I would look into a Negative Binomial because it has a variance larger than the mean. 

\end{document}